\vspace{0.5em}\noindent\textbf{HỆ GỢI Ý THẤT BẠI: DỮ LIỆU, KHÔNG PHẢI THUẬT
TOÁN}\par

Khi xây dựng hệ thống gợi ý sản phẩm với Amazon Personalize trong kỳ
thực tập tại FCAJ, chúng tôi phát hiện ra rằng hiệu suất kém thường xuất
phát từ chất lượng dữ liệu, không phải do chọn sai thuật toán.

\vspace{0.5em}\noindent\textbf{3 Bài Học Quan Trọng}\par

\vspace{0.5em}\noindent\textbf{1. Hiểu Cơ Chế Cốt Lõi Của Thuật
Toán}\par

\begin{itemize}
\tightlist
\item
  Amazon Personalize sử dụng \textbf{collaborative filtering} --- tìm
  những người dùng có hành vi tương tự và gợi ý dựa trên các mẫu hành vi
  đó.
\item
  \textbf{Điểm mấu chốt:} Nếu tương tác của người dùng là ngẫu nhiên, sẽ
  không có mẫu nào để thuật toán học.
\end{itemize}

\subsubsection{2. Chất Lượng Dữ Liệu \textgreater{} Chọn Thuật
Toán}

\begin{itemize}
\tightlist
\item
  Đổi recipe mà không sửa lỗi dữ liệu nền chỉ gây lãng phí thời gian và
  tiền bạc.
\item
  \textbf{Tập trung vào cấu trúc dữ liệu:}

  \begin{itemize}
  \tightlist
  \item
    Tạo các phân khúc người dùng thực tế (ví dụ: yêu thích điện tử, thời
    trang, sách)
  \item
    Mô phỏng hành vi thực tế (phiên truy cập, phễu chuyển đổi, phân phối
    lũy thừa)
  \item
    Thuật toán chỉ có thể học được những gì dữ liệu của bạn phản ánh
  \end{itemize}
\end{itemize}

\vspace{0.5em}\noindent\textbf{3. Dữ Liệu Mô Phỏng Phải Phản Ánh Thực
Tế}\par

\begin{itemize}
\tightlist
\item
  \texttt{random.choice()} tuy nhanh nhưng vô dụng cho việc kiểm thử có
  ý nghĩa.
\item
  \textbf{Thiết kế dữ liệu mô phỏng cẩn thận:}

  \begin{itemize}
  \tightlist
  \item
    Mô hình hóa phiên người dùng với các sản phẩm liên quan
  \item
    Áp dụng tỷ lệ chuyển đổi (xem\(\rightarrow\)giỏ,
    giỏ\(\rightarrow\)mua)
  \item
    Xây dựng các phân khúc người dùng với sở thích rõ ràng
  \end{itemize}
\end{itemize}

\vspace{0.5em}\noindent\textbf{So Sánh Kết Quả: Cùng Thuật Toán, Dữ Liệu Tốt
Hơn}\par

{\def\LTcaptype{none} % do not increment counter
\begingroup
\small
\setlength{\tabcolsep}{3pt}
\renewcommand{\arraystretch}{1.15}
\begin{longtable}[]{@{}
  >{\raggedright\arraybackslash}p{(\linewidth - 6\tabcolsep) * \real{0.2239}}
  >{\raggedright\arraybackslash}p{(\linewidth - 6\tabcolsep) * \real{0.2687}}
  >{\raggedright\arraybackslash}p{(\linewidth - 6\tabcolsep) * \real{0.2836}}
  >{\raggedright\arraybackslash}p{(\linewidth - 6\tabcolsep) * \real{0.2239}}@{}}
\toprule\noalign{}
\begin{minipage}[b]{\linewidth}\raggedright
Chỉ số
\end{minipage} & \begin{minipage}[b]{\linewidth}\raggedright
Dữ liệu ngẫu nhiên
\end{minipage} & \begin{minipage}[b]{\linewidth}\raggedright
Dữ liệu có cấu trúc
\end{minipage} & \begin{minipage}[b]{\linewidth}\raggedright
Cải thiện
\end{minipage} \\
\midrule\noalign{}
\endhead
\bottomrule\noalign{}
\endlastfoot
\textbf{Precision@5} & 0.0889 & 0.4348 & \textbf{Gấp 4.9 lần} \\
\textbf{NDCG@10} & 0.1799 & 0.6512 & \textbf{Gấp 3.6 lần} \\
\textbf{MRR@25} & 0.1216 & 0.7130 & \textbf{Gấp 5.9 lần} \\
\textbf{Coverage} & 0.8218 & 0.9505 & \textbf{+15.7\%} \\
\end{longtable}
\endgroup
}

\vspace{0.5em}\noindent\textbf{Những Điểm Rút
Ra:}\par

\begin{itemize}
\tightlist
\item
  MRR tăng gần 6 lần --- gợi ý phù hợp nay xuất hiện ở đầu danh sách nơi
  người dùng thực sự nhìn thấy
\item
  Coverage cải thiện 15.7\% --- hệ thống gợi ý đa dạng hơn, không chỉ
  gợi ý đi gợi ý lại các sản phẩm bán chạy
\item
  Một bộ dữ liệu có cấu trúc tốt có thể nhân hiệu suất lên 5-6 lần mà
  không cần thay đổi thuật toán
\end{itemize}

\vspace{0.5em}\noindent\textbf{Phân Tích Thành
Phần}\par

\textbf{Chuẩn Bị Dữ Liệu:} Nền tảng quan trọng nhất --- bao gồm phân
khúc người dùng, mô hình hành vi và các mẫu tương tác thực tế.

\textbf{Amazon Personalize:} Dịch vụ managed với các recipe có sẵn (AWS
cung cấp thuật toán, bạn cung cấp dữ liệu).

\textbf{Chỉ Số Đánh Giá:} Precision@5, NDCG@10, MRR@25, Coverage --- đo
chất lượng và đa dạng của gợi ý.

\textbf{Dữ Liệu Có Cấu Trúc:} Yêu cầu các phân khúc người dùng (điện tử,
thời trang, sách), tương tác theo phiên, phễu chuyển đổi và phân phối
lũy thừa.

\vspace{0.5em}\noindent\textbf{Điểm Mấu Chốt}\par

Công cụ và thuật toán chỉ chiếm 20\%. 80\% còn lại là \textbf{Chất Lượng
Dữ Liệu} và \textbf{Hiểu Rõ Lĩnh Vực Bài Toán} của bạn.

Khi các chỉ số thấp bất thường, \textbf{hãy kiểm tra dữ liệu trước khi
thay đổi thuật toán}.

\begin{center}\rule{0.5\linewidth}{0.5pt}\end{center}

\textbf{Bài viết gốc:}
\href{https://www.facebook.com/groups/awsstudygroupfcj/permalink/2229279517837008/?rdid=7zgI0zX0X33fQod5\#}{Facebook
AWS Study Group FCJ}
